\section{Explicit non-power-of-two witnesses and a dyadic lift bound}
\label{sec:nonpower-fff}

The computational nonexistence result at order $10$ does not extend to all even
non-powers of two.  The following finite claims are checked by direct
three-view cycle scans, not by an uncompleted solver search.

\begin{computationaltheorem}[Explicit non-power-of-two witnesses]
\label{cthm:nonpower-witnesses}
There are FFF Latin squares of orders $12$, $14$, and $36$.  The displayed
order-$12$ table is reduced.  Thus the former conjecture that FFF squares
of order greater than one exist only at powers of two is false.
\end{computationaltheorem}

\begin{table}[ht]
\centering
\begingroup\scriptsize\setlength{\tabcolsep}{3pt}
\begin{tabular}{r|*{12}{r}}
$L_{12}$ & 0&1&2&3&4&5&6&7&8&9&10&11 \\ \hline
0  & 0&1&2&3&4&5&6&7&8&9&10&11 \\
1  & 1&0&3&2&5&4&7&6&10&11&8&9 \\
2  & 2&3&0&1&7&6&5&4&9&8&11&10 \\
3  & 3&2&1&0&6&7&4&5&11&10&9&8 \\
4  & 4&7&5&6&8&11&9&10&3&0&2&1 \\
5  & 5&6&4&7&9&10&8&11&2&1&3&0 \\
6  & 6&5&7&4&10&9&11&8&1&2&0&3 \\
7  & 7&4&6&5&11&8&10&9&0&3&1&2 \\
8  & 8&9&10&11&3&2&1&0&6&7&4&5 \\
9  & 9&8&11&10&0&1&2&3&4&5&6&7 \\
10 & 10&11&8&9&2&3&0&1&5&4&7&6 \\
11 & 11&10&9&8&1&0&3&2&7&6&5&4 \\
\end{tabular}
\endgroup
\caption{An explicit reduced order-$12$ FFF witness.  Every row and column
is a permutation; all $198$ induced line-pair permutations across the three
views have only even cycles.}
\label{tab:order12-witness}
\end{table}

For odd $m>1$, let $h(m)$ be the least $a\geq0$ for which an FFF square of
order $m2^a$ exists.  The odd-order obstruction, the order-$6$ computation,
and Computational Theorem~\ref{cthm:nonpower-witnesses} give
$h(3)=2$ and $h(7)=1$; the order-$36$ table gives $h(9)\leq2$.
Nothing here decides existence at order $18$.

The classical polygon-and-centre round-robin factorisation gives a separate,
particularly simple family. Cheng and Sgueglia~\cite[Introduction]{ChengSgueglia2026Perfect}
recall this factorisation and its prime-order perfectness; the following
three-view FFF criterion is proved directly, without claiming priority for
the construction or an exhaustive literature classification.

\begin{theorem}[Prime round-robin family]
\label{thm:prime-round-robin-fff}
Let $p$ be an odd prime and $X=\mathbb F_p\cup\{\infty\}$. Define $M$ by
$M(\infty,\infty)=\infty$,
$M(\infty,x)=M(x,\infty)=x$ for finite $x$,
$M(x,x)=\infty$ for finite $x$, and
$M(x,y)=(x+y)/2$ for distinct finite $x,y$.
Then $M$ is Latin and is FFF if and only if $p\equiv3\pmod 8$.
\end{theorem}

\begin{proof}
In each finite row $a$, the affine map $y\mapsto(a+y)/2$ takes the value
$a$ only at $y=a$. Replacing that entry with $\infty$ and placing $a$ at
column $\infty$ makes the row a permutation. The infinity row is the
identity. Symmetry gives the Latin property for columns as well.

Write $R_a$ for a row permutation and $d=\operatorname{ord}_p(2)$.
$R_\infty$ is the identity, while $R_a$ swaps $a$ and $\infty$ and acts,
after translating $a$ to zero, as multiplication by $1/2$ on
$\mathbb F_p^*$. Hence an infinity/finite row pair has type
$2\,d^{(p-1)/d}$. For distinct finite $a,b$, put $P=R_b^{-1}R_a$.
Apart from $a$, $2b-a$, and $\infty$, the map is translation
$x\mapsto x+a-b$. At the exceptional inputs,
\[
 P(a)=b,\qquad P(2b-a)=\infty,\qquad P(\infty)=2a-b,
\]
and $P(b)=a$. In the single $p$-cycle of the translation, the segment
$2b-a\to b\to a\to2a-b$ is replaced by the $2$-cycle $(a,b)$ and
the connection $2b-a\to\infty\to2a-b$. The remaining cycle has length
$p-1$, including when $p=3$. Thus all finite/finite row pairs have type
$(2,p-1)$, and the row view is F exactly when $d$ is even. Symmetry gives
the identical column condition.

For a symbol $s$, let $T_s$ send a column to the row containing $s$ in it.
Then $T_\infty$ is the identity; for finite $s$,
\[
 T_s(\infty)=s,\qquad T_s(s)=\infty,\qquad
 T_s(x)=2s-x\quad(x\in\mathbb F_p\setminus\{s\}).
\]
Thus each $T_s$ is a fixed-point-free involution. An infinity/finite
symbol pair has type $2^{(p+1)/2}$. Affine relabelling of the finite
points reduces any two distinct finite symbols to $0,1$.
The relative permutation $T_1T_0$ has precisely the cycles
\[
 (0,1,3,5,\ldots,p-2),\qquad
 (\infty,2,4,6,\ldots,p-1),
\]
both of length $(p+1)/2$. Hence the symbol view is F exactly when
$p\equiv3\pmod4$.

For $p\equiv7\pmod8$, the supplement $(2/p)=1$ and Euler's criterion
make $d$ divide the odd number $(p-1)/2$, so the row and column views
are T. For $p\equiv3\pmod8$, $(2/p)=-1$ makes $d$ even; the symbol view
is also F. Primes $p\equiv1\pmod4$ already fail in the symbol view.
This proves the equivalence.
\end{proof}

By Dirichlet's theorem, infinitely many primes satisfy $p\equiv3\pmod8$;
apart from $p=3$, their orders $p+1$ are not powers of two. At $p=19$ the
table has order $20$. Together with the odd-order obstruction
Proposition~\ref{prop:odd-order} and the separately audited computational
exclusion of order $10$, this gives the
sharp threshold $h(5)=2$. This last equality depends on the exact
order-$10$ computation, whereas Theorem~\ref{thm:prime-round-robin-fff}
is algebraic.

\par\smallskip
\noindent\begin{minipage}{\linewidth}
\textbf{A second class at order $20$.}
The classical no-hexagon STS(19), with explicit starters in
Erskine--Griggs~\cite[Section 3.2]{ErskineGriggs2024CycleSwitch} and
classification context in Colbourn et al.~\cite{ColbournEtAl2010STS19}, yields a
Steiner loop whose 570 two-line permutations have only even cycles.
Its three-view cycle profile has no $18$-cycle, whereas the midpoint
table at $p=19$ has type $(2,18)$ in every row pair. Since the combined
cycle-type profile is invariant under paratopy, these two examples are
not in the same main class. This initial two-class control does not
reproduce the external STS classification. The separate trade-family
audit in Section~\ref{sec:fibred-growth} strengthens the represented
class lower bound to $514$.
\end{minipage}
\par\smallskip

\begin{theorem}[Three-point lift bound]
\label{thm:three-point-bound}
For every odd integer $m>1$, there is an FFF Latin square of order
$m2^{\operatorname{ord}_m(2)}$.  Consequently
\[
 h(m)\leq\operatorname{ord}_m(2),
\]
and FFF squares exist at every order $m2^a$ with
$a\geq\operatorname{ord}_m(2)$.
\end{theorem}

\begin{proof}
The case $m=3$ is supplied by the order-$12$ table above, since
$\operatorname{ord}_3(2)=2$.  Suppose $m>3$.  Set
$e=\operatorname{ord}_m(2)$, take $K=\mathbb F_{2^e}$, and choose
$r\in K^\times$ of multiplicative order $m$.  Its existence follows from
$m\mid(2^e-1)$ and cyclicity of $K^\times$.

For coefficients $s_i\in K$, $i\in\mathbb Z_m$, define a multiplication on
$\mathbb Z_m\times K$ by
\begin{equation}\label{eq:three-point-lift}
 (i,a)*(j,b)=(i+j,\ a+r^i b+s_i r^j).
\end{equation}
The quotient coordinate and the nonzero coefficient $r^i$ make this a Latin
square.  Put $t_i=r^{-i}s_i$ and regard $t_i$ as a function $t(r^i)$ on
$H=\langle r\rangle$.  For a coset $C$ of any nontrivial subgroup of $H$,
write
\[
 S_C=\sum_{x\in C}x\,t(x),\qquad
 T_C=\sum_{x\in C}x^{-1}t(x).
\]
We first verify that pairwise distinct $t_i$ and nonzero $S_C,T_C$ for all
such cosets imply FFF in all three views.

Two distinct lines in the same quotient fiber have relative permutations
that are nonzero translations of $K$: for rows the translation is
$r^{-i}(a+a')$, for columns it is $r^i(b+b')$, and for symbols it is
$r^{-(k-j)}(u+u')$ in each column fiber.  Each is an involution without
fixed points.  For distinct quotient rows $i,k$, put $d=i-k$ and let
$\ell$ be the order of $d$ in $\mathbb Z_m$.  Directly composing their row
maps gives
\[
 R_{(k,a')}^{-1}R_{(i,a)}(j,b)
 =\bigl(j+d,\ r^d b+r^{-k}(a+a')
       +r^{j-k}(s_i+r^d s_k)\bigr).
\]
The quotient shifts by $d$; its $\ell$-step fiber return is
\[
 b\longmapsto b+r^{j-i}(s_i+r^d s_k)
              =b+r^j(t_i+t_k).
\]
For distinct quotient columns $j,l$, the quotient shift is $j-l$ and the
one-step map, with $d=j-l$, is
\[
 C_{(l,b')}^{-1}C_{(j,b)}(i,a)
 =\bigl(i+d,\ a+r^i(b+b')+r^j s_i+r^l s_{i+d}\bigr).
\]
Its return on the associated coset $C$ is
\[
 a\longmapsto a+(r^j+r^l)S_C.
\]
For distinct quotient symbols $(k,u),(l,v)$, use the symbol-line map
\[
 S_{(k,u)}(j,b)=(k-j,\ u+r^{k-j}b+s_{k-j}r^j).
\]
With $d=l-k$, one step of $S_{(l,v)}^{-1}S_{(k,u)}$ is
\[
 (j,b)\longmapsto
 \bigl(j+d,\ b+r^{j-k}(u+v)+(1+r^d)r^j t_{k-j}\bigr).
\]
Summing over a quotient cycle gives the return
\[
 b\longmapsto b+(1+r^{l-k})r^k T_C.
\]
In each calculation the line-label terms vanish because a nontrivial
quotient cycle has odd length $\ell>1$ and the relevant geometric sum of
an $\ell$th root other than one is zero.  Every displayed multiplier outside
$S_C$ or $T_C$ is nonzero.  A nonzero characteristic-two translation has
order two, so each lifted cycle has even length $2\ell$.  This establishes
the asserted three-view criterion.

It remains to choose the coefficients.  Put $D=r+1$, $u=D^{-1}$, and start
with the injective function $t_0(1)=1$ and
$t_0(x)=(x+1)^{-1}$ for $x\in H\setminus\{1\}$.  Its values at
$1,r,r^2$ are $1,u,u^2$, which are distinct for $m>3$.  Cyclically permute
these three values:
\[
 t(1)=u,\qquad t(r)=u^2,\qquad t(r^2)=1,
 \qquad t(x)=(x+1)^{-1}\quad(x\notin\{1,r,r^2\}),
\]
and set $s_i=r^i t(r^i)$.  Injectivity is retained.

Let $C=aJ$ be a coset of a subgroup $J$ of odd order $\ell>1$, and write
$c=a^\ell$.  If $1\notin C$, then
$\prod_{x\in C}(X+x)=X^\ell+c$ and logarithmic differentiation at $X=1$
gives $\sum_{x\in C}(1+x)^{-1}=(1+c)^{-1}$.  Also
$\sum_{x\in C}x^{-1}=0$.  Any coset untouched by the patch therefore has
\[
 S_C=\frac{c}{1+c},\qquad T_C=\frac1{1+c},
\]
both nonzero.

Now let $J$ be proper, set $z=r^\ell$, and let
$A=(\ell-1)/2\pmod2$, interpreted as an element of $K$.  Its three
affected cosets $J,rJ,r^2J$ are distinct.  The value changes relative to
$t_0$ are $u+1$, $r/D^2$, and $r^2/D^2$, respectively.  Pairing
$x$ with $x^{-1}$ on $J\setminus\{1\}$ gives
$\sum_{x\in J\setminus\{1\}}(1+x)^{-1}=A$.  Including
$t_0(1)=1$, the unpatched moments on $J$ are $S_J(t_0)=A+1$ and
$T_J(t_0)=A$.  The other two cosets use the preceding unpatched formula.
Adding the appropriate value changes yields
\[
\begin{array}{c|cc}
 C&S_C&T_C\\ \hline
 J&A+u&A+1+u\\[2pt]
 rJ&\dfrac{z+r^2}{(1+z)D^2}&\dfrac{z+r^2}{(1+z)D^2}\\[8pt]
 r^2J&\dfrac{[z(r^2+r+1)+r^2]^2}{(1+z^2)D^2}
      &\dfrac{z^2+r^2}{(1+z^2)D^2}
\end{array}
\]
All denominators are nonzero.  The $J$ entries are nonzero because
$u\notin\{0,1\}$.  The $rJ$ numerator could vanish only if
$\ell\equiv2\pmod m$, and the last $T$ numerator only if
$\ell\equiv1\pmod m$; neither is possible for a proper divisor
$1<\ell<m$.  The only possible failure is therefore
\begin{equation}\label{eq:bad-three-point-coset}
 r^\ell(r^2+r+1)=r^2.
\end{equation}
At most one exponent modulo $m$ satisfies this equation.  If it holds for
a proper divisor $\ell$, replace $r$ by $r^{-1}$.  A bad proper divisor
$j$ for the inverse would require
$r^2+r+1=r^j=r^{2-\ell}$, so
$j=m+2-\ell>m/2$.  Every proper divisor of odd $m$ is at most $m/3$;
there is no such $j$.  Thus either $r$ or its inverse makes every proper
coset moment nonzero.

For the full-group coset $H$, all three patched points lie together.  With
$A=(m-1)/2\pmod2$, direct summation gives
\[
 S_H=1+A+r^2+\frac r{r+1},\qquad
 T_H=A+\frac r{r+1}.
\]
When $m\equiv3\pmod4$, these are
$r(r^2+r+1)/(r+1)$ and $1/(r+1)$; the first could vanish only if $r$
had order three.  When $m\equiv1\pmod4$, they are
$(r^3+r^2+1)/(r+1)$ and $r/(r+1)$.  The cubic is irreducible over
$\mathbb F_2$ and its roots have order seven, incompatible with
$m\equiv1\pmod4$.  Hence both moments are nonzero.  The three-view
criterion proves that~\eqref{eq:three-point-lift} is FFF.

Finally, multiplying by the table of an elementary abelian $2$-group
preserves FFF by Corollary~\ref{cor:fff-product}, giving every exponent
$a\geq e$.
\end{proof}

\begin{remark}[The exceptional affine base case]
At $m=3$, over the prescribed minimal field
$K=\mathbb F_{2^{\operatorname{ord}_3(2)}}=\mathbb F_4$,
the affine formula~\eqref{eq:three-point-lift} cannot be FFF,
although Table~\ref{tab:order12-witness} gives an order-$12$ FFF square.
Put $t_i=r^{-i}s_i$. Row-F requires distinct $t_i$; column-F and symbol-F
require $S=t_0+rt_1+r^2t_2\ne0$ and $T=t_0+r^2t_1+rt_2\ne0$.
Common translation and nonzero scaling give $(t_0,t_1)=(0,1)$.
Since $r^2+r+1=0$, $t_2=r$ forces $T=0$, whereas $t_2=r^2$
forces $S=0$.
This argument concerns $\mathbb F_4$, not the same formula over all larger
binary fields. As a scope control, over
$\mathbb F_{16}=\mathbb F_2[x]/(x^4+x+1)$ take $r=x^2+x$,
$(t_0,t_1,t_2)=(0,1,x)$, and $s_i=r^it_i$.
Direct three-view checks give an FFF square of order $48$.
This is not a new order-$48$ existence claim: products of the
order-$12$ witness already supply that order.
\end{remark}

The finite checks in the associated companion run evaluate all subgroup
coset moments for $m=5,7,9,15,21$ and all line-pair cycles at orders
$56$, $80$, and $240$.  They control the implementation; the universal
quantifier in Theorem~\ref{thm:three-point-bound} rests on the argument
above.  The theorem does not determine $h(m)$ in general, resolve order
$18$, classify higher-order main classes, or establish literature priority.
